% debye_boltzmann_ic_walkthrough.tex
%
% A pedagogical walkthrough of the Debye--Boltzmann initial condition
% used by the PNP--BV forward solver. Written for a reader with strong
% applied-math / numerical background but limited electrochemistry
% intuition.

\documentclass[11pt]{article}
\usepackage[margin=1in]{geometry}
\usepackage{amsmath,amssymb,bm}
\usepackage{xcolor}
\usepackage[most]{tcolorbox}
\usepackage{microtype}
\usepackage{parskip}
\usepackage{siunitx}
\usepackage{enumitem}
\usepackage{hyperref}
\usepackage{booktabs}

\setlist[itemize]{leftmargin=*,topsep=2pt,itemsep=3pt}
\setlist[enumerate]{leftmargin=*,topsep=2pt,itemsep=3pt}

\hypersetup{colorlinks=true, linkcolor=blue!50!black, urlcolor=blue!50!black}

\newcommand{\VT}{V_{T}}
\newcommand{\Eeq}{E_{\mathrm{eq}}}
\newcommand{\VRHE}{V_{\mathrm{RHE}}}
\newcommand{\lD}{\lambda_D}

\definecolor{boxgray}{gray}{0.95}
\tcbset{
  notebox/.style={
    enhanced, breakable,
    colback=boxgray, colframe=black!40,
    boxrule=0.4pt, arc=2pt,
    left=8pt, right=8pt, top=6pt, bottom=6pt,
    fonttitle=\bfseries,
  }
}

\title{The Debye--Boltzmann Initial Condition\\[2pt]
{\large A walkthrough for the electrochemistry-light reader}}
\author{Jake Weinstein \\ \small PNP--BV forward solver, May 2026}
\date{}

\begin{document}
\maketitle

\begin{abstract}
\noindent
This note explains, from electrochemistry first principles, what the
\texttt{debye\_boltzmann} initial condition does, why it exists, and
where each formula comes from. It is meant to be read end-to-end
without prior PNP/electrochemistry background. Equations are given
where they help, but the goal is to build the picture, not to be
exhaustive. Code anchor:
\texttt{Forward/bv\_solver/forms\_logc.py:581--952}.
\end{abstract}

\tableofcontents
\vspace{1em}\hrule\vspace{1em}

%% =====================================================================
\section{What problem is the IC actually solving?}

The forward solver is a Firedrake finite-element discretization of a
coupled nonlinear PDE system:
\begin{itemize}
\item One Nernst--Planck (NP) equation per dynamic species
      ($\mathrm{O_2}$, $\mathrm{H_2O_2}$, $\mathrm{H^+}$),
\item One Poisson equation for the electrostatic potential $\varphi$,
\item Two Butler--Volmer (BV) reaction-rate boundary conditions on the
      electrode surface (R$_1$: $\mathrm{O_2}\to\mathrm{H_2O_2}$,
      R$_2$: $\mathrm{H_2O_2}\to\mathrm{H_2O}$).
\end{itemize}

Steady-state is reached by pseudo-time-stepping with backward Euler
plus Newton's method (SNES). Newton needs a starting guess. That
starting guess is the \emph{initial condition} (``IC'') discussed
here.

\begin{tcolorbox}[notebox, title={A point of vocabulary}]
``IC'' here is overloaded with the time-evolution sense of initial
condition. We do not care about transient dynamics --- the only thing
we want from the IC is for it to be \emph{close enough to the
steady-state solution that Newton converges}. A bad IC means Newton
diverges or stalls. A good IC means Newton converges in a handful of
iterations. That is the entire point.
\end{tcolorbox}

\subsection{Why a simple IC fails}

The naive IC --- ``set every concentration to its bulk value, set
$\varphi$ to a linear ramp from $\varphi_{\mathrm{applied}}$ at the
electrode to $0$ in the bulk'' --- works at low applied voltage but
fails badly at moderate to high voltage. To see why, you need to
understand what happens physically near the electrode.

%% =====================================================================
\section{Two-minute electrochemistry primer}

\subsection{The setup}

Picture a 1D domain. Coordinate $y \in [0, 1]$ (nondimensionalized).

\begin{itemize}
\item $y = 0$ is the \textbf{electrode} (a metal disk you applied a
      voltage to).
\item $y = 1$ is the \textbf{bulk} (well-mixed electrolyte far from
      the electrode).
\item In between is the electrolyte: water with dissolved
      $\mathrm{H^+}$, $\mathrm{ClO_4^-}$ (counterion),
      $\mathrm{O_2}$, and $\mathrm{H_2O_2}$.
\end{itemize}

When you hold the electrode at potential $\varphi_{\mathrm{applied}}$,
two things happen at once:
\begin{enumerate}
\item Charges \textbf{rearrange} to screen the field --- positive ions
      are pushed away from a positive electrode, negative ions are
      pulled in (and vice-versa). This is purely electrostatic and
      happens essentially instantaneously.
\item BV \textbf{reactions} consume $\mathrm{O_2}$ and $\mathrm{H^+}$
      and produce $\mathrm{H_2O_2}$, then consume $\mathrm{H_2O_2}$
      and $\mathrm{H^+}$ to make $\mathrm{H_2O}$. This sets up
      diffusion gradients away from the electrode.
\end{enumerate}

These two processes happen on \emph{wildly different length scales}.

\subsection{The Debye length: where screening lives}

The screening length is the \emph{Debye length}:
\begin{equation}
  \lD \;=\; \sqrt{\frac{\varepsilon \VT}{2 F^2 c_b}},
\end{equation}
where $\varepsilon$ is the electrolyte permittivity, $\VT = RT/F
\approx 25.7$~mV at room temperature, $F$ is the Faraday constant, and
$c_b$ is the bulk ion concentration.

For pH 4 with a 0.1 M counterion, $\lD \sim 1$~nm. The full domain
might be \SI{100}{\micro\meter} thick. So the screening region is
$10^{-4}$ of the domain. The PNP equations are stiff in exactly the
sense that interesting physics is concentrated in a tiny corner.

\subsection{The Boltzmann distribution: charge follows potential}

Inside the Debye layer, ions reach a local equilibrium: their
electrochemical potential becomes constant. For species $i$ with
charge $z_i$,
\begin{equation}
  \mu_i \;=\; \ln c_i \,+\, z_i \varphi \;=\; \mathrm{const}
  \;\;\Longrightarrow\;\;
  \boxed{\,c_i(y) \;=\; c_i^{\mathrm{bulk}}\,\exp\bigl(-z_i\,\psi(y)\bigr)\,}
\end{equation}
where $\psi(y) = \varphi(y) - \varphi^{\mathrm{bulk}}$ is the
deviation from bulk potential, in units of $\VT$. (All $\varphi$
quantities in this document are nondimensionalized by $\VT$, so
factors of $\VT$ vanish.)

This is the \textbf{Boltzmann distribution}. It is the equilibrium
profile a single charged species takes in an electric field. \emph{The
entire IC machinery is built around this formula.}

\paragraph{Numerical concern.} If the electrode is held at
$\psi_D \equiv \varphi_{\mathrm{electrode}} - \varphi^{\mathrm{bulk}}
= +30$, then for the proton ($z_H = +1$),
\begin{equation*}
  c_H^{\mathrm{surface}}
   \;=\; c_H^{\mathrm{bulk}} \cdot e^{-30}
   \;\approx\; c_H^{\mathrm{bulk}} \cdot 10^{-13}.
\end{equation*}
Thirteen orders of magnitude. A Newton solver starting from
$c_H = c_H^{\mathrm{bulk}}$ everywhere has to discover this depletion
in one step, which it cannot.

\subsection{Outer vs.\ inner: two regions, two physics}

Combining the two effects above gives a clean two-region picture:

\begin{tcolorbox}[notebox, title={The two-region picture}]
\textbf{Outer region} ($y \gtrsim \lD$, the bulk of the domain):
mass transport dominates. Concentrations are nearly linear in $y$
(steady-state diffusion under a flux boundary condition is just a
linear profile). Charge density is negligible
($\sum_i z_i c_i \approx 0$, electroneutrality).
\medskip

\textbf{Inner region} ($y \lesssim \lD$, the Debye layer): screening
dominates. Concentrations follow Boltzmann. Charge density is
nonzero, and that's exactly what Poisson's equation balances.
\medskip

\textbf{Outer surface}: the conceptual boundary at $y \approx \lD$
where the two regions meet. Quantities with subscript $o$ (e.g.
$H_o$, $\varphi_o$) are evaluated here.
\end{tcolorbox}

The IC strategy is to build each region separately and glue them
together at the outer surface.

%% =====================================================================
\section{Step 1 --- Outer linear profiles}

In the outer region, only diffusion matters. At steady state with a
constant flux at the electrode and Dirichlet (bulk) at $y = 1$, the
profile is linear. The species fluxes equal the BV reaction rates:
\begin{align}
  J_{\mathrm{O_2}}     &=  R_1                  & \text{(O$_2$ consumed by R$_1$)} \\
  J_{\mathrm{H_2O_2}}  &= -R_1 + R_2            & \text{(produced by R$_1$, consumed by R$_2$)} \\
  J_{\mathrm{H^+}}     &=  R_1 + R_2            & \text{(2 protons consumed in each reaction; the 2 cancels, see below)}
\end{align}

The corresponding linear profiles, written in terms of bulk values
($O_b$, $P_b$, $H_b$) and outer-surface values ($O_s$, $P_s$, $H_o$),
are
\begin{equation}
  \boxed{\;
  O_s = O_b - \frac{R_1}{D_O},\qquad
  P_s = P_b + \frac{R_1 - R_2}{D_P},\qquad
  H_o = H_b - \frac{R_1 + R_2}{D_H}.\;}
\end{equation}

\begin{tcolorbox}[notebox, title={Why the proton denominator is $D_H$, not $2 D_H$}]
The two-proton stoichiometry of each reaction would naively give
$(R_1+R_2)/(2 D_H) \cdot 2$ in the numerator. But \emph{ambipolar
transport} doubles the effective proton diffusivity: because
electroneutrality enforces that $\mathrm{H^+}$ and $\mathrm{ClO_4^-}$
move together, the slower partner sets the pace, and the field
drags the faster one. The result is that the effective proton
diffusivity is $2 D_H$, not $D_H$, and the factor of 2 cancels
against the stoichiometric 2. The bare $D_H$ in the denominator is
correct. (Comment in
\texttt{forms\_logc.py:798--800} explicitly warns: ``do not
correct''.)
\end{tcolorbox}

So the outer-region values $(O_s, P_s, H_o)$ are determined as soon as
we know the rates $(R_1, R_2)$. But the rates depend on the surface
concentrations, which depend on $\psi_D$, which depends on $H_o$\dots

\section{Step 2 --- The Picard loop}

We have a self-referential algebraic system:
\begin{itemize}
\item Rates $R_1, R_2$ depend on the \emph{surface} (inner) value
      $H_s = H_o \exp(-\psi_D)$.
\item $H_o$ depends on $R_1, R_2$ via the outer linear profile.
\item $\psi_D$ depends on $\varphi_o$, which depends on $H_o$.
\end{itemize}
This is solved by Picard iteration. The whole loop is closed-form
algebra --- no PDE solve, no Newton. Each iteration is cheap.

\subsection{The 2$\times$2 system}

Inside the loop, given a current guess for $(R_1, R_2)$:
\begin{enumerate}
\item Compute outer values:
\[
  H_o = H_b - (R_1 + R_2)/D_H,
\]
and from outer-region electroneutrality plus the Boltzmann profile
of the counterion,
\[
  \varphi_o = \ln(H_o / c_{\mathrm{ClO_4}}^{\mathrm{bulk}}),
  \qquad
  \psi_D = \varphi_{\mathrm{applied}} - \varphi_o.
\]
\item Compute the surface proton concentration:
\[
  H_s = H_o \exp(-\psi_D).
\]
\item Build BV rate constants. R$_1$ has two branches (cathodic
``forward'', anodic ``backward''); R$_2$ here keeps only the
cathodic branch:
\begin{align*}
  A_1 &= k_1 \,(H_s/H_{\mathrm{ref}})^2\, \exp(-\alpha_1 n_e \eta_1)
        & \text{R$_1$ cathodic,}\\
  B_1 &= k_1 \,\exp\bigl((1-\alpha_1)\, n_e \eta_1\bigr)
        & \text{R$_1$ anodic,}\\
  A_2 &= k_2 \,(H_s/H_{\mathrm{ref}})^2\, \exp(-\alpha_2 n_e \eta_2)
        & \text{R$_2$ cathodic.}
\end{align*}
The $H_s$ enters squared because each reaction consumes 2 protons.
\item Solve the 2$\times$2 system for new rates:
\begin{equation}
  \begin{bmatrix}
    1 + A_1/D_O + B_1/D_P & -B_1/D_P\\
    -A_2/D_P              & 1 + A_2/D_P
  \end{bmatrix}
  \begin{bmatrix} R_1 \\ R_2 \end{bmatrix}
  =
  \begin{bmatrix} A_1 O_b - B_1 P_b \\ A_2 P_b \end{bmatrix}.
\end{equation}
\item Damped update: $R_j \leftarrow (1-\omega) R_j^{\mathrm{old}} +
\omega R_j^{\mathrm{new}}$ with $\omega = 0.5$.
\item Loop until $|\Delta R_1| + |\Delta R_2| < 10^{-6}$, max 50
iterations.
\end{enumerate}

\subsection{Where this matrix came from}

The matrix is just a rearrangement of mass balance + the linear-profile
formulas. Take R$_1$:
\[
  R_1 \;=\; A_1\,O_s \,-\, B_1 P_s
       \;=\; A_1\,(O_b - R_1/D_O) \,-\, B_1\,(P_b + (R_1-R_2)/D_P)
\]
Collecting $R_1, R_2$ on the left:
\[
  \bigl(1 + A_1/D_O + B_1/D_P\bigr) R_1 \,-\, (B_1/D_P) R_2
   \;=\; A_1 O_b - B_1 P_b.
\]
Same exercise on R$_2$ gives the second row. So this matrix is just
``write down what each rate equals, substitute the linear profile,
collect terms''. Nothing fancy.

\subsection{Failure handling}

If the Picard loop diverges, hits NaN, or hits a singular Jacobian,
the IC silently \textbf{falls back to a trivial linear-$\varphi$ IC}.
This happens in practice below $\VRHE = +0.20$~V, where the H$^+$
mass-transport limit makes the loop oscillate. The fallback is
recorded in
\texttt{ctx[`initializer\_fallback'] = True}.

%% =====================================================================
\section{Step 3 --- Inner profile: the Debye layer}

After Picard, we have $(R_1, R_2, H_o, \varphi_o, \psi_D, O_s, P_s)$
--- all scalar values at the outer surface. We now need to extend
each field $u_i(y)$ and $\varphi(y)$ over the \emph{whole mesh}, not
just point values.

The outer region is easy --- linear profiles already given. The hard
part is the inner Debye layer, where $\psi$ goes from
$\psi_D$ at the electrode to roughly $0$ at the outer surface, over a
distance of $\lD$.

\subsection{The Gouy--Chapman profile (linear-Debye limit)}

For a 1:1 symmetric electrolyte, the nonlinear Poisson--Boltzmann
equation in 1D admits a closed-form solution:
\begin{equation}
  \psi_{\mathrm{GC}}(y)
   \;=\; 4\,\mathrm{atanh}\!\left[\tanh(\psi_D / 4)\,
                                  e^{-y/\lD}\right].
\end{equation}
This is the \emph{Gouy--Chapman profile}. At small $\psi_D$ (the
linear-Debye limit) it collapses to a pure exponential
$\psi(y) \approx \psi_D\,e^{-y/\lD}$. For larger $\psi_D$ it stays
finite but the inner ramp is steeper.

\paragraph{Numerical rewrite.} Direct evaluation of $4\,\mathrm{atanh}$
overflows when $\psi_D \gtrsim 30$. Use the identity
\[
  4\,\mathrm{atanh}(x) \;=\; 2\ln\!\left(\frac{1+x}{1-x}\right),
\]
which stays well-behaved as $|x| \to 1$ (just clamp the argument
slightly inside $\pm 1$). The implementation does exactly this.

\subsection{When linear-Debye is not enough: Bikerman steric}

Pure Boltzmann assumes ions are point charges. At the electrode,
$\psi_D = +30$ would give a counterion concentration of
$c_b \cdot e^{30} \approx c_b \cdot 10^{13}$ --- which is physically
absurd because the ions have a finite volume and can only pack so
tightly.

\textbf{Bikerman steric} is the simplest fix: penalize the local
packing fraction in the chemical potential. Ions can pack up to
$1/a$ where $a$ is the (nondimensional) ion volume. Mathematically,
the modified Boltzmann distribution becomes
\begin{equation}
  c_b(y) \;=\;
     c_b^{\mathrm{bulk}}
     \frac{\bigl(1 - \sum_j a_j c_j(y)\bigr)\,
           \exp(-z_b \psi(y))}
          {\theta_b + a_b c_b^{\mathrm{bulk}}\exp(-z_b\psi(y))},
\end{equation}
which caps $c_b \to 1/a_b$ as $\psi \to -\infty$. This denominator
is what bounds the Poisson source above
$\VRHE = +0.66$~V where the unbounded Boltzmann co-ion would have
made the Poisson residual blow up.

\subsection{The BKSA composite-$\psi$ profile}

Once steric effects matter, the $\psi(y)$ profile changes character.
Instead of one smooth tanh-like curve, it has \emph{two zones}:

\begin{tcolorbox}[notebox, title={Two-zone Debye structure under steric saturation}]
\textbf{Saturation zone} (close to electrode): the counterion is at
close-packing density $c \approx 1/a$. Charge density is roughly
constant in this zone. Poisson's equation
$-\varepsilon \nabla^2 \psi = \rho$ with constant $\rho$ gives a
\emph{linear} $\psi$ profile (think of it like a parallel-plate
capacitor field):
\[
  \psi(y) \;=\; \mathrm{sgn}(\psi_D)\,(|\psi_D| - \alpha y),
  \qquad 0 \le y < y_{\mathrm{match}}.
\]

\textbf{Outer Debye zone} (linear-Debye regime): once enough screening
has happened that $|\psi|$ drops below the saturation threshold
$\psi_{\mathrm{sat}}$, ions are no longer at close-packing and we
return to ordinary exponential decay:
\[
  \psi(y) \;=\; \mathrm{sgn}(\psi_D)\,\psi_{\mathrm{sat}}\,
                \exp\!\bigl(-(y - y_{\mathrm{match}})/\lD\bigr),
  \qquad y \ge y_{\mathrm{match}}.
\]
\end{tcolorbox}

The match point $y_{\mathrm{match}}$ and slope $\alpha$ come from
matched asymptotics (Bazant--Kilic--Storey--Ajdari, BKSA):
\begin{align}
  \nu &\;=\; 2\,a_{\mathrm{Cl}}\,c_{\mathrm{Cl}}^{\mathrm{bulk}}
        & \text{(steric crowding parameter)} \\
  \psi_{\mathrm{sat}} &\;=\; \ln(2/\nu)
        & \text{(potential at which counterion close-packs)} \\
  \alpha &\;=\; \sqrt{\frac{2}{\nu \lD^2}\,
                       \ln\!\bigl(1 + \nu(\cosh\psi_D - 1)\bigr)}
        & \text{(slope in the saturated zone)} \\
  y_{\mathrm{match}} &\;=\; (|\psi_D| - \psi_{\mathrm{sat}})/\alpha
        & \text{(thickness of the saturated zone)}
\end{align}

If $|\psi_D| < \psi_{\mathrm{sat}}$, the saturated zone has zero
thickness and we fall back to the pure linear-Debye exponential
(or the Gouy--Chapman tanh, depending on regime).

\paragraph{Where this comes from.} The first integral of the
steric Poisson--Boltzmann equation is
$\frac{1}{2} \lD^2 (\psi')^2 = (1/\nu) \ln(1 + \nu(\cosh\psi - 1))$.
Setting $\psi = \psi_{\mathrm{sat}}$ gives $\alpha$ as the slope at
the saturation boundary, and the linear extrapolation backward to
$\psi = \pm \psi_D$ gives $y_{\mathrm{match}}$. The exponential outer
piece is the linearization of Gouy--Chapman in the small-$\psi$
limit. So the composite profile is just \emph{matched asymptotics
glued together at the boundary where saturation ends}.

%% =====================================================================
\section{Step 4 --- The multispecies $\gamma$ correction}

Even with the right $\psi(y)$ in hand, naively writing
$c_i(y) = c_i^{\mathrm{outer}}(y) \cdot \exp(-z_i \psi(y))$
\textbf{violates} the Bikerman packing constraint at the surface --- the
constraint and the simple Boltzmann formula are inconsistent.

The fix is a closed-form multiplicative correction $\gamma(y)$ that
restores consistency. From multispecies zero-flux equilibrium with
the steric chemical potential, one derives
\begin{equation}
  \boxed{\;
  \gamma(y) \;=\;
     \frac{1}{1 + \sum_j a_j\,c_j^{\mathrm{anchor}}(y)\,
                    \bigl(\exp(-z_j \psi(y)) - 1\bigr)}.
  \;}
\end{equation}
Two convenient features of this expression:
\begin{itemize}
\item \textbf{Neutral species drop out} automatically, because
      $e^{0} - 1 = 0$. Only charged species ($\mathrm{H^+}$ and
      $\mathrm{ClO_4^-}$) contribute to the denominator.
\item In the linear-Debye limit ($\psi$ small), $e^{-z\psi} - 1
      \approx -z \psi$ and $\gamma \to 1$ smoothly. So the correction
      vanishes in the regime where it shouldn't matter.
\end{itemize}

The IC then seeds every species on the saturated steric manifold:
\begin{align}
  u_{\mathrm{O_2}}^{\mathrm{IC}}(y)    &= \ln O_{\mathrm{outer}}(y)  + \ln \gamma(y), \\
  u_{\mathrm{H_2O_2}}^{\mathrm{IC}}(y) &= \ln P_{\mathrm{outer}}(y)  + \ln \gamma(y), \\
  u_{\mathrm{H^+}}^{\mathrm{IC}}(y)    &= \ln H_{\mathrm{outer}}(y) - \psi(y) + \ln \gamma(y),\\
  \varphi^{\mathrm{IC}}(y)             &= \ln\!\bigl(H_{\mathrm{outer}}(y) / c_{\mathrm{ClO_4}}^{\mathrm{bulk}}\bigr) + \psi(y).
\end{align}

\begin{tcolorbox}[notebox, title={What the pieces are doing, in plain English}]
\begin{itemize}
\item ``$\ln$ outer'' : the bulk-to-outer-surface diffusion ramp.
\item ``$-\psi(y)$'' on H$^+$ : Boltzmann depletion of the proton in
      the Debye layer (a $+1$ ion is repelled from a positive
      electrode, so its concentration drops).
\item ``$+\psi(y)$'' on $\varphi$ : the actual electrostatic
      potential profile, going from the bulk-electroneutrality
      reference up to $\varphi_{\mathrm{applied}}$ at the electrode.
\item ``$+\ln\gamma$'' on every species : the steric correction
      that makes the Boltzmann profile obey the packing constraint.
\end{itemize}
\end{tcolorbox}

%% =====================================================================
\section{Step 5 --- Why $\mu_H$ is the Right primary variable}

The production solver does \emph{not} carry $u_H = \ln c_H$ as the
proton's primary unknown. It carries the
\textbf{electrochemical potential}
\[
  \mu_H \;=\; u_H + z_H \varphi \;=\; u_H + \varphi
  \qquad (z_H = +1).
\]
Concentrations are reconstructed as $c_H = \exp(\mu_H - \varphi)$.

\subsection{Why this helps}

Inside the Debye layer where H$^+$ is in local equilibrium,
$\nabla \mu_H \approx 0$. So $\mu_H$ is \emph{nearly constant in $y$
across the Debye layer} --- exactly where $u_H$ and $\varphi$ are
each varying by tens of log-units. Newton sees a smooth unknown
in the regime where the alternatives are wildest.

\subsection{Exact cancellation under the Debye--Boltzmann IC}

A clean correctness check: under the IC of Section 4 (ignoring
$\gamma$ for clarity),
\begin{align*}
  \mu_H^{\mathrm{IC}}(y)
  &= u_H^{\mathrm{IC}}(y) + \varphi^{\mathrm{IC}}(y) \\
  &= \bigl[\ln H_{\mathrm{outer}}(y) - \psi(y)\bigr]
   + \bigl[\ln\bigl(H_{\mathrm{outer}}(y)/c_{\mathrm{Cl}}^{\mathrm{bulk}}\bigr) + \psi(y)\bigr]\\
  &= 2 \ln H_{\mathrm{outer}}(y) - \ln c_{\mathrm{Cl}}^{\mathrm{bulk}}.
\end{align*}
The diffuse-layer $\psi$ --- the wildest-varying piece --- cancels
exactly. The IC for $\mu_H$ contains \emph{only} the slowly-varying
outer profile of H$^+$ and a constant offset. That is as smooth as
it gets.

%% =====================================================================
\section{Putting it all together}

\subsection{Algorithm summary}

\begin{enumerate}
\item Read bulk values, diffusivities, BV constants, and steric sizes
      from \texttt{solver\_params}.
\item Picard-iterate the algebraic 2$\times$2 system to get
      $(R_1, R_2)$, $(O_s, P_s, H_o)$, and $\psi_D$ at the outer
      surface. Bail to linear-$\varphi$ on failure.
\item Build the \textbf{outer linear profiles} $O_{\mathrm{outer}}(y)$,
      $P_{\mathrm{outer}}(y)$, $H_{\mathrm{outer}}(y)$ as Firedrake
      expressions on the mesh.
\item Build the \textbf{inner $\psi(y)$ profile}: pure Gouy--Chapman
      if the electrolyte is not steric, BKSA composite if it is.
\item Build the \textbf{multispecies $\gamma(y)$ correction} (when
      steric).
\item Interpolate $u_i^{\mathrm{IC}}(y)$ and $\varphi^{\mathrm{IC}}(y)$
      onto the mesh; copy into \texttt{ctx[`U']} and
      \texttt{ctx[`U\_prev']}.
\item Hand to Newton.
\end{enumerate}

\subsection{What it buys you, in numbers}

\begin{center}
\begin{tabular}{lll}
\toprule
\textbf{Metric} & \textbf{Linear-$\varphi$ IC} & \textbf{Debye--Boltzmann IC} \\
\midrule
Newton iters at $\VRHE = +0.30$~V & $\sim 50$ & $\sim 10$--30
   \;($1.7$--$6.3\times$ fewer) \\
Cold-start ceiling                & $+0.30$~V & $+0.60$~V \\
Warm-walk reach (with Stern)      & $+0.60$~V & $+1.00$~V (15/15) \\
\bottomrule
\end{tabular}
\end{center}

The Bikerman-consistent IC alone (composite-$\psi$ + multispecies-$\gamma$,
``Option 2b'') accounts for the +0.4 V gain from $+0.60$ to $+1.00$ V on
the production stack.

\subsection{Soft spots to be aware of}

\begin{itemize}
\item \textbf{Picard can fail silently.} The fallback to linear-$\varphi$
      is automatic but is a no-op below $\VRHE = +0.20$~V (the
      H$^+$ mass-transport limit makes the outer loop oscillate).
\item \textbf{The IC is not on the adjoint tape.} The whole body
      runs under \texttt{firedrake.adjoint.stop\_annotating()}.
      The IC is an initial \emph{guess}, not a control --- you don't
      want pyadjoint to think the Picard iterations are part of the
      gradient computation.
\item \textbf{It is still a guess.} If the production-stack residual
      changes (e.g.\ a different counterion model, a different BV
      formulation), the IC can become inconsistent with what Newton
      is converging to, and convergence degrades. The IC and the
      residual must agree about what the saturated steric manifold
      looks like; this is hard rule \#3 in
      \texttt{CLAUDE.md}.
\end{itemize}

%% =====================================================================
\section{One-paragraph explainer for the meeting}

\begin{tcolorbox}[notebox, title={The version to say out loud}]
The Debye--Boltzmann IC is an analytical sketch of the steady-state
solution that we hand to Newton as a starting guess. It builds the
solution in two pieces: an outer linear-diffusion profile from bulk
to a fictitious ``outer surface'' just outside the Debye layer, and
an inner Boltzmann-style screening profile from the outer surface
down to the electrode. The two are stitched together using a
self-consistent algebraic 2$\times$2 Picard solve for the BV reaction
rates, a Gouy--Chapman or BKSA composite $\psi(y)$ profile inside the
Debye layer, and a multispecies activity correction $\gamma(y)$ that
keeps the Bikerman packing constraint satisfied. With the IC,
Newton converges 1.7--6.3$\times$ faster, and unblocks cold starts
that the linear IC could not reach. It is not part of the residual
and not on the adjoint tape --- it is a guess, optimized to be
right almost everywhere.
\end{tcolorbox}

\end{document}
